% Source: https://artofproblemsolving.com/community/c3h3181929_2021_abmc_online_december_contest__actonboxborough_math_competition
% Stable identity: aops-topic-3181929-q11
Let departure times in hours after noon be $u\in[0,3]$ and $v\in[1,4]$. If their absolute separation is $\delta<3$, the earlier driver meets the later while returning; the meeting position is the fraction $1-\dfrac{\delta}{3}$ of the route from Pittsburgh. It is closer to Pittsburgh exactly when $\dfrac{3}{2}<\delta<3$. In the 9-area departure rectangle, the region $v-u>\dfrac{3}{2}$ has area $\dfrac{1}{2}(\dfrac{5}{2})^2=\dfrac{25}{8}$; remove the $v-u\ge3$ triangle of area $\dfrac{1}{2}$, leaving $\dfrac{21}{8}$. The opposite region $u-v>\dfrac{3}{2}$ has area $\dfrac{1}{2}(\dfrac{1}{2})^2=\dfrac{1}{8}$, and automatically has separation below 3. Total favourable area is $\dfrac{11}{4}$, giving probability $\dfrac{11}{36}$. The requested sum is $\boxed{47}$.
\begin{asy}
usepackage("amsmath");
size(7cm);
draw((0,1)--(3,1)--(3,4)--(0,4)--cycle);filldraw((0,1.5)--(0,3)--(1,4)--(2.5,4)--cycle,lightgray);filldraw((2.5,1)--(3,1)--(3,1.5)--cycle,lightgray);label("$u$",(3,1),SE);label("$v$",(0,4),NW);label("departure rectangle",(1.5,.5));
\end{asy}
